\section{Introduction}
Under distribution shift, evaluation depends not only on the metrics reported but also on the checkpoint selected for reporting.
Prior work shows that model-selection rules and stopping criteria can materially change out-of-distribution conclusions \citep{gulrajani2020domainbed,cha2021swad,arpit2022ensembleaverages,yang2023changeishard,teterwak2025ermpp,naganuma2025pretrainedmodelselection,berta2025refinethencalibrate}.
We study this as an evaluation-practice failure mode: when a shaped objective induces its own validation proxy, what reliability claim does the proxy-selected report support under shift?

Many modern pipelines optimize a shaped objective rather than plain cross-entropy or raw MSE \citep{lin2017focalloss,reed2014bootstrapping,zhang2018generalizedcrossentropy}.
We focus on one subset, which we call \emph{suppressive}: objectives or target transforms that reduce gradient pressure on high-loss examples.
Soft clipping is a controlled classification example, and winsorized targets in heavy-tailed regression are a natural regression example.
Once chosen, such an objective induces its own validation score, which we call the \emph{proxy}.
The resulting workflow is simple: train with the shaped objective, pick the checkpoint with the best validation proxy, and report headline held-out metrics.

Our question is whether that workflow remains trustworthy under shift.
For suppressive objectives, it often does not.
The proxy-favored checkpoint can look best on the proxy, and in some cases even on held-out accuracy, while doing worse on held-out loss, calibration, and group- or tail-sensitive reliability.

\paragraph{Scope.}
Our claim is specific.
We do not claim that validation-loss selection is a universal selector, or that every shaped objective creates the same failure.
Most evidence is a baseline-referenced workflow comparison: the objective, trajectory, and selector change together.
Selector-only conclusions are limited to explicitly labeled same-trajectory selector and admissibility checks.
The largest mean accuracy-up / reliability-down pattern appears on Camelyon17 finetuning, where proxy-only selection raises mean held-out accuracy while worsening held-out loss, validation-side calibration, and a teacher-tail stress diagnostic.
The teacher-tail diagnostic is a fixed held-out bank of examples that a reference model finds hard; we use it as a stress check, not as a semantic or clinical subgroup.
The primary reliability evidence is held-out loss, calibration, and, where available, official group-sensitive metrics.

The expected part is that suppressing high-loss examples can hurt tail-sensitive metrics.
The evaluation hazard is that scalarized checkpoint selection can hide that damage and legitimize the checkpoint that caused it.

The core anchors are Camelyon17 ERM, Camelyon17 finetuning, and frozen CivilComments.
CelebA and ACSIncome serve narrower scope checks: CelebA shows transience, and ACSIncome tests the regression analogue.
Additional GCE, hard-bootstrap, teacher-free tail, averaging, objective-family, and time-course analyses test mechanism and scope rather than expanding the main claim.

Prior work already established that selector choice matters under shift \citep{gulrajani2020domainbed,cha2021swad,arpit2022ensembleaverages,yang2023changeishard,teterwak2025ermpp,naganuma2025pretrainedmodelselection,berta2025refinethencalibrate}.
Our narrower contribution is to show a specific reporting failure: suppressive proxy-selected workflows can report checkpoints that look attractive on the proxy or on accuracy while failing standard held-out reliability checks.
Within-trajectory selector disagreement is the sharpest evidence, and baseline-referenced comparisons show what a suppressive pipeline can present as a reported win.
The rest of the paper tests when this hazard persists, when a validation-loss diagnostic helps, and whether the risky regime is tied to suppression rather than to objective shaping broadly.

\paragraph{Contributions.}
As an evaluation-design contribution, the paper makes three claims.
\begin{itemize}[leftmargin=*,nosep]
\item \textbf{Reporting failure.} Checkpoint selection is part of the evaluation protocol: under objectives that downweight high-loss examples, proxy-selected reporting can make less reliable checkpoints look successful under shift.
\item \textbf{Scope separation.} We separate workflow comparisons, where objective, trajectory, and selector change together, from same-trajectory selector/admissibility checks. The main evidence uses held-out loss, calibration, and official WILDS worst-group loss, with teacher-tail used as supporting fixed-bank stress evidence.
\item \textbf{Reporting rule, not optimizer.} Report both proxy-selected and validation-loss-selected checkpoints, and treat harmful disagreement on held-out reliability metrics as an unvalidated reliability claim.
\end{itemize}
